% ---
% Hallazgos de implementación y lecciones técnicas aprendidas (OE-6).
% Sección 6.7 del reporte técnico.
% Entradas que lo respaldan: H-012, H-013, H-014, H-019, H-032, H-034, D-002, D-014, D-015, D-017, D-031.
% ---

\subsection{Hallazgos de implementación}
\label{sec:lecciones_implementacion}

La materialización práctica del pipeline experimental de aprendizaje automático cuántico en simulación clásica
condujo al descubrimiento de un conjunto de limitaciones, ambigüedades de especificación y particularidades numéricas
tanto del formalismo cuántico como de la biblioteca de desarrollo. Estos hallazgos constituyen una contribución técnica
directa a la respuesta del objetivo específico OE-6, al documentar los factores que determinan la reproducibilidad y
la viabilidad operacional del cómputo cuántico simulado. Para cada caso se distingue explícitamente entre propiedades
fundamentales del formalismo cuántico y particularidades de implementación de la biblioteca Qiskit.

% fuente: H-012, D-002; B1_benchmark_complejidad.csv
\textbf{Identidad estructural entre circuitos en mapas de Pauli.} Al comparar la fidelidad cuántica entre los
estados generados por el mapa \texttt{zz\_feature\_map} (condición C4) y el mapa general \texttt{pauli\_feature\_map}
configurado con $\text{paulis} = [\text{'Z'},\text{'ZZ'}]$ (especificación inicial de C5), se obtuvo un valor
de fidelidad idéntico a $1.000000000000$ tanto para $k = 4$ como para $k = 8$. De manera concurrente, las 18 filas del análisis
de complejidad de circuitos en B1 resultaron exactamente idénticas en profundidad cruda (50 a 112), profundidad transpilada
(58 a 126), compuertas CX (63 a 285) y conteo total de compuertas (187 a 725). Este comportamiento obedece a un
\textbf{detalle de implementación de Qiskit}: en la biblioteca, la función \texttt{zz\_feature\_map} está definida de
forma literal como una invocación a \texttt{pauli\_feature\_map} con el conjunto $[\text{'Z'},\text{'ZZ'}]$. Por tanto,
ambas opciones representaban el mismo circuito y no constituían condiciones experimentales independientes. Ante esta evidencia,
se descartó esa especificación y se evaluaron configuraciones alternativas con idéntico coste de compuertas
entrelazadoras ($n_{cx} = 12$) promediando sobre 200 entradas aleatorias en $[0, \pi]^4$.

% fuente: H-034, D-015, D-031
\textbf{Inoperancia del operador Pauli X tras la capa de Hadamard.} Al formular la condición C5 mediante
$\text{paulis} = [\text{'X'},\text{'ZZ'}]$ con el propósito de alternar el eje de rotación de primer orden,
la fidelidad respecto a un circuito que omite por completo los términos lineales en $X$ (un mapa puramente $ZZ$) arrojó
un valor mínimo de $0.99999999999999$ para 20 entradas al azar con $reps = 1$, y sobre las lesiones reales de masas la
discrepancia máxima en las matrices de kernel fue de $5.9 \times 10^{-15}$. Este fenómeno corresponde a una
\textbf{propiedad del formalismo cuántico}: la capa inicial de compuertas Hadamard prepara a cada qubit en el estado
$|+\rangle$, el cual es autoestado de la compuerta de Pauli $X$ con autovalor $+1$. La rotación $R_X(2x_i) = H P(2x_i) H$
actúa sobre su propio autoestado multiplicando al qubit únicamente por una fase compleja escalar. Sobre el estado producto
global, estas fases individuales se agregan en una fase global indetectable que no afecta ninguna medición ni producto
interno. Como consecuencia, se redefinió la condición
C5 utilizando $\text{paulis} = [\text{'Y'},\text{'ZZ'}]$, donde la rotación en $Y$ no preserva el autoestado ecuatorial
y modifica sustantivamente la codificación previa al entrelazamiento $ZZ$.

% fuente: H-013, D-014
\textbf{Nulidad idéntica del vector de expectativas sin capa variacional.} Al evaluar el vector de valores
esperados $\langle Z_i \rangle$ directamente a la salida del mapa de características $ZZ$ prescindiendo del \textit{ansatz}
variacional, el resultado fue $\langle Z_i \rangle = 0$ de forma exacta para todo qubit y cualquier vector de entrada $x$.
Este resultado responde a una \textbf{propiedad del formalismo cuántico}: partiendo del estado de superposición uniforme
$H^{\otimes k}|0\rangle$ (en el cual todas las amplitudes poseen módulo $2^{-k/2}$), todas las compuertas que componen el
mapa $ZZ$ (rotaciones de fase $R_Z$ y bloques $CX-R_Z-CX$) son operadores diagonales en la base computacional. Las matrices
diagonales multiplican cada amplitud por factores de fase compleja pero preservan de forma estricta sus magnitudes originales.
Por lo tanto, la probabilidad marginal de medir cero satisface $P(\text{qubit}_i = 0) = 1/2$, lo que impone $\langle Z_i \rangle = 0$.
Toda la información inyectada por el mapa reside en las fases relativas, respecto a las cuales la medición en $Z$ es ciega.
En respuesta, se concluyó que bajo la Arquitectura~B el \textit{ansatz} variacional es metodológicamente indispensable:
sus parámetros $\boldsymbol\theta$ deben fijarse con una semilla declarada (semilla 42) para que las rotaciones $R_Y(\boldsymbol\theta)$
proyecten las diferencias de fase a poblaciones medibles. Asimismo, se clarificó que el kernel de fidelidad no sufre esta
limitación al operar sobre los vectores de estado completos en el espacio de Hilbert.

% fuente: H-019, H-014, H-016; B3_origen_suelo_h016.csv, B3_shots_muestreo_real.csv
\textbf{Valores por defecto en estimadores y modelado de disparos en Aer.} En las evaluaciones iniciales
de sensibilidad al número de disparos (*shots*), el error frente a la supuesta referencia analítica se estancaba en un suelo
constante de $\approx 0.013$ hasta un millón de disparos, y el error parecía insensible entre 512 y 8\,192 disparos.
Este comportamiento combinó una \textbf{propiedad del formalismo cuántico} con dos \textbf{detalles de implementación de Qiskit}.
En el formalismo, la dispersión por disparos decae como $1/\sqrt{n}$ con varianza $(1 - \langle Z_i \rangle^2)/n$. En la biblioteca,
la clase \texttt{EstimatorQNN} se construye con una precisión por defecto de $0.015625$ ($= 1/\sqrt{4096}$) y sobrescribe cualquier
precisión fijada en el estimador subyacente; al transferir dicho valor a \texttt{StatevectorEstimator}, este último suma ruido
gaussiano $\mathcal{N}(0, \text{precisión})$ incluso en simulación de vector de estado, generando un error medio analítico de
$0.015625 \sqrt{2/\pi} = 0.01247$ en la referencia teórica. Además, el estimador \texttt{EstimatorV2} de Aer no simula disparos
discretos, sino que calcula la expectativa analítica y le adiciona ruido gaussiano continuo, generando una dependencia circular
con pendiente log-log de $-0.5065$. En consecuencia, se invalidó aquel hallazgo preliminar y la sección de disparos de B1; para el Módulo 5 se fijó
que toda expectativa exacta debe evaluarse con \texttt{Statevector} directo o primitivas con precisión cero explícita (\texttt{default\_precision = 0.0});
y para modelar disparos reales se adoptó en B3 el muestreador \texttt{SamplerV2} de Aer a partir de cadenas de bits y, en el estudio
del OE-6, el muestreo multinomial sobre las probabilidades exactas $|\psi(\mathbf{x})|^2$, equivalente al anterior sin ruido de
compuerta. En B3 se verificaron pendientes de $-0.4986$ en dispersión y $-0.4983$ en MAE sin suelo espurio hasta $2^{20}$ disparos.

% fuente: H-032, D-017
\textbf{Aceleración analítica mediante productos internos matriciales.} La estimación inicial de tiempo de cálculo
en B1 proyectaba aproximadamente 41 minutos por matriz de kernel cuántico de $200 \times 200$ mediante \texttt{FidelityQuantumKernel}
(superando las 2.75~h para las cuatro matrices del protocolo), mientras que la implementación con \texttt{FidelityStatevectorKernel}
evaluó la misma matriz exacta en 1.5\,s. Este hallazgo constituye un \textbf{detalle de implementación del simulador clásico}:
\texttt{FidelityQuantumKernel} construye y ejecuta circuitos cuánticos independientes de tipo \textit{compute-uncompute} por cada
par de muestras ($O(N^2)$ ejecuciones de circuito), mientras que \texttt{FidelityStatevectorKernel} aprovecha la disponibilidad
simultánea de los $N$ vectores de estado en memoria clásica ($O(N)$ simulaciones) y calcula la matriz completa mediante productos
internos vectoriales directos. Se adoptó \texttt{FidelityStatevectorKernel} para agilizar los módulos M6 y derivados en más de
tres órdenes de magnitud ($\sim 1\,000\times$), explicitando en el reporte que dicho atajo computacional es una ventaja
exclusiva de la simulación clásica que no existe en procesadores cuánticos físicos.

% fuente: D-031
\textbf{Obsolescencia de \texttt{PauliFeatureMap} en Qiskit~2.x.} A partir de Qiskit~2.1, la clase
\texttt{PauliFeatureMap} está marcada como obsoleta en favor de la función \texttt{pauli\_feature\_map()}, que devuelve un
\texttt{QuantumCircuit} estándar. Se trata de un \textbf{detalle de implementación de la API de Qiskit}, sin consecuencia
sobre el circuito que se construye. El código del proyecto emplea las funciones \texttt{pauli\_feature\_map()} y
\texttt{zz\_feature\_map()}.
